\documentclass{article}

\usepackage{amssymb}
\textwidth 6.in
\textheight 8.5in
\hoffset-0.5in
\voffset-.5in


\title{Dynamica \\
A computer package for the study of Difference Equations}
\author{M. R. S. Kulenovic and Orlando Merino \\
Department of Mathematics \\
University of Rhode Island \\
Kingston, RI 02881}
\date{}

\begin{document}
\maketitle

{\it Dynamica} is a computer package, written in Mathematica,
intended for use in the study of difference equations and discrete dynamical systems.

Let $k$ be a positive integer and let
$f:R^k \rightarrow R$ be a given function.
A relation of the form
\begin{equation}
\label{difference equation}
x_{n+1} = f(x_n, x_{n-1} , \ldots , x_{n-k+1})
\end{equation}
is called a difference equation of order $k$.

Difference equations arise in virtually all areas of science,
mathematics and engineering.  Difference equations are discrete
analogs of differential equations.

{\it Dynamica} implements a series of tools and techniques
of algebraic, numerical, and graphical nature used
in the study of difference equations (\ref{difference equation}).
These include:
\begin{itemize}
\item  Finding equilibrium and periodic points.
\item  Classifying the stability character of equilibrium and periodic points.
\item  Semicycle analysis of solutions.
\item Calculation and visualization of invariants (first integrals).
\item Calculation and visualization of Lyapunov functions.
\item Plotting bifurcation diagrams.
\item Visualization of stable and unstable manifolds.
\item Calculation and visualization of Lyapunov numbers.
\item Calculation of Box Dimension.
\end{itemize}

{\it Dynamica} consists of a series of programs and a two sets of notebooks.

The first set of {\it Dynamica} notebooks serve as a tutorial of the 
programs and the different techniques.  The necessary concepts and 
results are provided in the text, as well as references to research publications, and suggested problems.  However, the presentation style 
emphasizes use of the software and not theoretical discussions.

The second set of {\it Dynamica} notebooks consists of case studies 
of well known difference equations, which come under name such as Henon, Lyness, Mira, Todd, and others.  An interesting feature of the notebooks
is that they show how to prove, with help of the tools provided with
{\it Dynamica}, many results published recently.  

{\it Dynamica} is easy to learn and use.  In very little time the new user
may be ready to study complex difference equations of which not
much (maybe almost nothing!) is known.
\end{document}
